% Options for packages loaded elsewhere
\PassOptionsToPackage{unicode}{hyperref}
\PassOptionsToPackage{hyphens}{url}
\PassOptionsToPackage{dvipsnames,svgnames,x11names}{xcolor}
\documentclass[
  french,
  10pt,
  a4paper,
]{article}
\usepackage{xcolor}
\usepackage[margin=2.2cm]{geometry}
\usepackage{amsmath,amssymb}
\setcounter{secnumdepth}{-\maxdimen} % remove section numbering
\usepackage{iftex}
\ifPDFTeX
  \usepackage[T1]{fontenc}
  \usepackage[utf8]{inputenc}
  \usepackage{textcomp} % provide euro and other symbols
\else % if luatex or xetex
  \usepackage{unicode-math} % this also loads fontspec
  \defaultfontfeatures{Scale=MatchLowercase}
  \defaultfontfeatures[\rmfamily]{Ligatures=TeX,Scale=1}
\fi
\usepackage{lmodern}
\ifPDFTeX\else
  % xetex/luatex font selection
\fi
% Use upquote if available, for straight quotes in verbatim environments
\IfFileExists{upquote.sty}{\usepackage{upquote}}{}
\IfFileExists{microtype.sty}{% use microtype if available
  \usepackage[]{microtype}
  \UseMicrotypeSet[protrusion]{basicmath} % disable protrusion for tt fonts
}{}
\makeatletter
\@ifundefined{KOMAClassName}{% if non-KOMA class
  \IfFileExists{parskip.sty}{%
    \usepackage{parskip}
  }{% else
    \setlength{\parindent}{0pt}
    \setlength{\parskip}{6pt plus 2pt minus 1pt}}
}{% if KOMA class
  \KOMAoptions{parskip=half}}
\makeatother
\usepackage{longtable,booktabs,array}
\usepackage{calc} % for calculating minipage widths
% Correct order of tables after \paragraph or \subparagraph
\usepackage{etoolbox}
\makeatletter
\patchcmd\longtable{\par}{\if@noskipsec\mbox{}\fi\par}{}{}
\makeatother
% Allow footnotes in longtable head/foot
\IfFileExists{footnotehyper.sty}{\usepackage{footnotehyper}}{\usepackage{footnote}}
\makesavenoteenv{longtable}
\ifLuaTeX
\usepackage[bidi=basic,shorthands=off,]{babel}
\else
\usepackage[bidi=default,shorthands=off,]{babel}
\fi
\ifLuaTeX
  \usepackage{selnolig} % disable illegal ligatures
\fi
\setlength{\emergencystretch}{3em} % prevent overfull lines
\providecommand{\tightlist}{%
  \setlength{\itemsep}{0pt}\setlength{\parskip}{0pt}}
% Fichier LaTeX généré par docs/build/build_pdf.sh à partir de 04_autres_fonctions_claude.md.
% Compilation (depuis le dossier de ce fichier), deux fois pour la table des matières :
%   pdflatex 06_Guide4_Autres_fonctions_Claude.tex
% Pour insérer une capture : déposez le fichier PNG au nom indiqué dans le cadre « CAPTURE À INSÉRER »
% (dossier ../images/...), puis recompilez. Vous pouvez aussi remplacer la commande \CredixCapture{...}{...}
% par votre propre \includegraphics.
\newcommand{\CredixShortTitle}{Guide 4 : autres fonctions}
% Préambule commun aux PDF du dossier de remise CREDIX (compilé avec pdfLaTeX via pandoc)
\usepackage[scaled=.92]{helvet}
\renewcommand{\familydefault}{\sfdefault}
\usepackage{pifont}
\usepackage{amssymb}
% Caractères Unicode utilisés dans les documents (pdfLaTeX ne les connaît pas d'office)
\DeclareUnicodeCharacter{03C1}{\ensuremath{\rho}}
\DeclareUnicodeCharacter{2265}{\ensuremath{\geq}}
\DeclareUnicodeCharacter{2264}{\ensuremath{\leq}}
\DeclareUnicodeCharacter{2248}{\ensuremath{\approx}}
\DeclareUnicodeCharacter{2192}{\ensuremath{\rightarrow}}
\DeclareUnicodeCharacter{2190}{\ensuremath{\leftarrow}}
\DeclareUnicodeCharacter{2212}{\ensuremath{-}}
\DeclareUnicodeCharacter{2713}{\ding{51}}
\DeclareUnicodeCharacter{2714}{\ding{52}}
\DeclareUnicodeCharacter{2717}{\ding{55}}
\DeclareUnicodeCharacter{25B2}{\ensuremath{\blacktriangle}}
\DeclareUnicodeCharacter{25CB}{\ensuremath{\circ}}
\DeclareUnicodeCharacter{2500}{-}
\DeclareUnicodeCharacter{2502}{|}
\DeclareUnicodeCharacter{250C}{+}
\DeclareUnicodeCharacter{2510}{+}
\DeclareUnicodeCharacter{2514}{+}
\DeclareUnicodeCharacter{2518}{+}
\DeclareUnicodeCharacter{251C}{+}
\DeclareUnicodeCharacter{2524}{+}
\DeclareUnicodeCharacter{252C}{+}
\DeclareUnicodeCharacter{2534}{+}
\DeclareUnicodeCharacter{253C}{+}
\DeclareUnicodeCharacter{0192}{\textit{f}}
\DeclareUnicodeCharacter{2011}{-}
\DeclareUnicodeCharacter{202F}{\,}
\DeclareUnicodeCharacter{2009}{\,}
\DeclareUnicodeCharacter{00D7}{\ensuremath{\times}}
\DeclareUnicodeCharacter{25A0}{\ensuremath{\blacksquare}}
\DeclareUnicodeCharacter{2022}{\ensuremath{\bullet}}
\DeclareUnicodeCharacter{2026}{\ldots}
\DeclareUnicodeCharacter{2460}{\ding{172}}
\DeclareUnicodeCharacter{2461}{\ding{173}}
\DeclareUnicodeCharacter{2462}{\ding{174}}
\DeclareUnicodeCharacter{2463}{\ding{175}}
\DeclareUnicodeCharacter{2464}{\ding{176}}
\DeclareUnicodeCharacter{2465}{\ding{177}}
\DeclareUnicodeCharacter{2466}{\ding{178}}
\DeclareUnicodeCharacter{2467}{\ding{179}}
\usepackage{fvextra}
\DefineVerbatimEnvironment{verbatim}{Verbatim}{breaklines=true,breakanywhere=true,fontsize=\footnotesize,frame=single,framerule=0.3pt,rulecolor=\color[gray]{0.7},framesep=2mm,xleftmargin=1mm,xrightmargin=1mm}
\fvset{breaklines=true,breakanywhere=true}
\usepackage{fancyhdr}
\usepackage{tcolorbox}
\usepackage{colortbl}
\usepackage{enumitem}
\definecolor{credixblue}{HTML}{1F3A5F}
\definecolor{credixgray}{HTML}{5F6B7A}
\definecolor{boxbg}{HTML}{F3F5F8}
\setlength{\emergencystretch}{4em}
\renewcommand{\arraystretch}{1.3}
\setlength{\parindent}{0pt}
\setlength{\parskip}{0.55em}
\setlist{itemsep=0.15em,topsep=0.3em}
\pagestyle{fancy}
\fancyhf{}
\fancyhead[L]{\small\color{credixgray}\textsc{\CredixShortTitle}}
\fancyhead[R]{\small\color{credixgray}CREDIX --- Dossier de remise}
\fancyfoot[C]{\small\color{credixgray}\thepage}
\renewcommand{\headrulewidth}{0.3pt}
\renewcommand{\footrulewidth}{0pt}
\fancypagestyle{plain}{\fancyhf{}\fancyfoot[C]{\small\color{credixgray}\thepage}\renewcommand{\headrulewidth}{0pt}}
\usepackage{titlesec}
\titleformat{\section}{\Large\bfseries\color{credixblue}}{}{0em}{}[\vspace{-0.2em}\color{credixblue}\titlerule]
\titleformat{\subsection}{\large\bfseries\color{credixblue}}{}{0em}{}
\titleformat{\subsubsection}{\normalsize\bfseries\color{credixblue}}{}{0em}{}
% Cadre de repère pour une capture manquante
\newcommand{\CaptureManquante}[2]{%
  \begin{tcolorbox}[colback=boxbg,colframe=credixgray,boxrule=0.5pt,arc=1mm,left=3mm,right=3mm,top=3mm,bottom=3mm,halign=center]
  \textbf{CAPTURE À INSÉRER}\\[0.4em]#1\\[0.4em]{\footnotesize\ttfamily #2}
  \end{tcolorbox}}

% Images : insérées si le fichier existe au moment de la compilation, sinon cadre « CAPTURE À INSÉRER ».
% Pour ajouter une capture : déposer le PNG au nom indiqué dans le cadre, puis recompiler.
\usepackage{graphicx}
\newcommand{\CredixCapture}[2]{%
  \IfFileExists{#2}{%
    \begin{center}\includegraphics[width=\linewidth,height=0.75\textheight,keepaspectratio]{#2}\end{center}}%
  {\CaptureManquante{#1}{\detokenize{#2}}}}
\newcommand{\CredixFigure}[1]{%
  \begin{center}\includegraphics[width=0.95\linewidth,height=0.8\textheight,keepaspectratio]{#1}\end{center}}
\usepackage{bookmark}
\IfFileExists{xurl.sty}{\usepackage{xurl}}{} % add URL line breaks if available
\urlstyle{same}
\hypersetup{
  pdftitle={Guide 4 : autres fonctions de Claude},
  pdfauthor={SINGHE PENKA Hendrix Donavan},
  pdflang={fr},
  colorlinks=true,
  linkcolor={credixblue},
  filecolor={Maroon},
  citecolor={Blue},
  urlcolor={credixblue},
  pdfcreator={LaTeX via pandoc}}

\author{}
\date{}

\begin{document}

\begin{titlepage}
\thispagestyle{empty}
\centering
\vspace*{3.2cm}
{\color{credixgray}\large\textsc{Projet de fin d'études --- ENSPY, Université de Yaoundé I}\par}
\vspace{2.2cm}
{\Huge\bfseries\color{credixblue} Guide 4 : autres fonctions de Claude\par}
\vspace{0.8cm}
{\Large Fichiers Excel, Word, PowerPoint, Skills, artefacts, connecteurs\par}
\vspace{2cm}
{\color{credixblue}\rule{0.5\linewidth}{0.8pt}\par}
\vspace{1.2cm}
{\large SINGHE PENKA Hendrix Donavan\par}
{\color{credixgray}Dépôt GitHub : HendrixPenka2/credix\par}
\vfill
{\color{credixgray}\small Version du 19/09/2026\par}
\end{titlepage}


\renewcommand*\contentsname{Table des matières}
{
\setcounter{tocdepth}{2}
\tableofcontents
}
\section{Guide 4 : les autres fonctions de
Claude}\label{guide-4--les-autres-fonctions-de-claude}

\textbf{Créer des fichiers Excel, Word, PowerPoint et PDF, transformer
un PDF en présentation, utiliser les Skills, les artefacts et les
connecteurs, et fabriquer soi-même les PDF de ce dossier.}

\begin{quote}
\textbf{À propos des captures d\textquotesingle écran.} Ce guide
contient des cadres « CAPTURE À INSÉRER », numérotés de 04-01 à 04-06.
Chaque cadre indique ce qu\textquotesingle il faut montrer et le nom du
fichier attendu, dans
\texttt{docs/\allowbreak{}images/\allowbreak{}guides/\allowbreak{}}. La
liste complète est dans
\texttt{docs/\allowbreak{}guides/\allowbreak{}CAPTURES\_\allowbreak{}A\_\allowbreak{}PRENDRE.\allowbreak{}md}.

\textbf{À propos des menus.} Les noms de menus de Claude changent avec
le temps, et leur libellé français peut différer de
l\textquotesingle anglais. Ceux de ce guide ont été \textbf{vérifiés sur
les pages officielles le 19 septembre 2026} (sources en section 11).

\textbf{Ce qui est vérifié, et ce qui ne l\textquotesingle est pas.} Les
fonctions décrites (formats, limites, menus) viennent de la
documentation officielle. Les \textbf{recettes} (sections 3 et 4) sont
des façons de les utiliser : \textbf{essayez-les une fois avant de les
présenter}, et relisez toujours le résultat.
\end{quote}

\begin{center}\rule{0.5\linewidth}{0.5pt}\end{center}

\subsection{1. À quoi sert ce guide}\label{1-uxe0-quoi-sert-ce-guide}

Les guides 1 à 3 montrent comment concevoir une application avec Claude.
Ce guide, \textbf{à part}, répond à une autre question : \textbf{que
sait faire Claude en dehors de la conception ?}

Vous y trouverez :

\begin{itemize}
\tightlist
\item
  la création de \textbf{fichiers} (Excel, Word, PowerPoint, PDF)
  directement dans une conversation ;
\item
  une recette complète pour \textbf{transformer un PDF en présentation
  PowerPoint}, par exemple un mémoire en diaporama de soutenance ;
\item
  les \textbf{Skills}, qui apprennent à Claude vos consignes de travail
  une fois pour toutes ;
\item
  les \textbf{artefacts} et les \textbf{connecteurs} (Google Drive,
  GitHub, etc.) ;
\item
  ce que \textbf{Claude Code} sait faire en plus
  d\textquotesingle écrire du code ;
\item
  la fabrication \textbf{sans intelligence artificielle} de PDF, de
  documents Word et de présentations à partir de fichiers Markdown,
  comme pour ce dossier.
\end{itemize}

\begin{center}\rule{0.5\linewidth}{0.5pt}\end{center}

\subsection{2. Créer des fichiers avec
Claude}\label{2-cruxe9er-des-fichiers-avec-claude}

\subsubsection{2.1 Ce que Claude peut
produire}\label{21-ce-que-claude-peut-produire}

Claude peut \textbf{créer et modifier des fichiers} directement dans une
conversation. Vous les téléchargez ensuite, ou vous les enregistrez dans
Google Drive.

\begin{longtable}[]{@{}
  >{\raggedright\arraybackslash}p{(\linewidth - 4\tabcolsep) * \real{0.2153}}
  >{\raggedright\arraybackslash}p{(\linewidth - 4\tabcolsep) * \real{0.1625}}
  >{\raggedright\arraybackslash}p{(\linewidth - 4\tabcolsep) * \real{0.5922}}@{}}
\toprule\noalign{}
\begin{minipage}[b]{\linewidth}\raggedright
Type de fichier
\end{minipage} & \begin{minipage}[b]{\linewidth}\raggedright
Extension
\end{minipage} & \begin{minipage}[b]{\linewidth}\raggedright
Exemples d\textquotesingle usage
\end{minipage} \\
\midrule\noalign{}
\endhead
\bottomrule\noalign{}
\endlastfoot
\textbf{Tableur Excel} & \texttt{.xlsx} & Un tableau de données trié, un
budget, une analyse de résultats \\
\textbf{Présentation PowerPoint} & \texttt{.pptx} & Un diaporama de
soutenance, un support de formation \\
\textbf{Document Word} & \texttt{.docx} & Un rapport, une lettre, un
compte rendu \\
\textbf{PDF} & \texttt{.pdf} & Un document figé prêt à envoyer \\
\textbf{Image de graphique} & \texttt{.png} & Une courbe, un histogramme
tracés à partir de vos données \\
\textbf{Script Python} & \texttt{.py} & Un programme
d\textquotesingle analyse que vous pouvez relancer \\
\end{longtable}

Cette fonction est disponible sur \textbf{tous les abonnements}
(Gratuit, Pro, Max, Team, Enterprise), sur le site web,
l\textquotesingle application de bureau et l\textquotesingle application
mobile.

\subsubsection{2.2 L\textquotesingle activer}\label{22-lactiver}

\begin{enumerate}
\def\labelenumi{\arabic{enumi}.}
\tightlist
\item
  Cliquez sur votre profil (en bas à gauche), puis \textbf{Paramètres}
  (\emph{Settings}).
\item
  Ouvrez \textbf{Capacités} (\emph{Capabilities}).
\item
  Activez l\textquotesingle option \textbf{« Code execution and file
  creation »} (exécution de code et création de fichiers).
\end{enumerate}

\CredixCapture{Paramètres, onglet Capacités : l'option de création de fichiers}{../images/guides/04-01_parametres_capacites.png}

\emph{Figure 1. L\textquotesingle option « Code execution and file
creation ».}

Sur les abonnements d\textquotesingle équipe (Team, Enterprise),
l\textquotesingle option est activée par défaut au niveau de
l\textquotesingle organisation, et un propriétaire peut la désactiver.

\subsubsection{2.3 Les limites à
connaître}\label{23-les-limites-uxe0-connauxeetre}

\begin{longtable}[]{@{}
  >{\raggedright\arraybackslash}p{(\linewidth - 2\tabcolsep) * \real{0.4667}}
  >{\raggedright\arraybackslash}p{(\linewidth - 2\tabcolsep) * \real{0.5033}}@{}}
\toprule\noalign{}
\begin{minipage}[b]{\linewidth}\raggedright
Limite
\end{minipage} & \begin{minipage}[b]{\linewidth}\raggedright
Valeur
\end{minipage} \\
\midrule\noalign{}
\endhead
\bottomrule\noalign{}
\endlastfoot
Taille d\textquotesingle un fichier créé ou téléchargé & \textbf{30 Mo}
au maximum \\
Fichiers envoyés dans une conversation & \textbf{20 fichiers} au maximum
par conversation \\
Taille d\textquotesingle un fichier envoyé dans une conversation & 500
Mo au maximum ; \textbf{30 Mo} dans un projet \\
PDF envoyé & \textbf{1 000 pages} au maximum \\
Ce que Claude lit dans un PDF de \textbf{100 pages ou moins} & Le
\textbf{texte et les éléments visuels} (images, graphiques) \\
Ce que Claude lit dans un PDF de \textbf{101 à 1 000 pages} & Le
\textbf{texte seulement} \\
Formats de documents acceptés & PDF, DOCX, CSV, TXT, HTML, ODT, RTF,
EPUB, JSON, XLSX \\
Formats d\textquotesingle images acceptés & JPEG, PNG, GIF, WebP \\
\end{longtable}

\begin{quote}
\textbf{Conséquence pratique.} Un mémoire de plus de 100 pages est lu
\textbf{sans ses graphiques}. Si les schémas comptent, envoyez les
chapitres séparément (moins de 100 pages chacun), ou ajoutez les images
vous-même à la fin.
\end{quote}

\subsubsection{2.4 Un point de
sécurité}\label{24-un-point-de-suxe9curituxe9}

Quand Claude peut créer des fichiers, il exécute du code dans un
environnement séparé. Deux précautions vous concernent :

\begin{itemize}
\tightlist
\item
  \textbf{Le partage public} d\textquotesingle une conversation qui
  contient des fichiers créés est \textbf{désactivé} sur les abonnements
  Gratuit, Pro et Max.
\item
  \textbf{L\textquotesingle accès à Internet} de cet environnement est
  réglable. Une page web piégée pourrait, en théorie, pousser Claude à
  envoyer vos données ailleurs : \textbf{n\textquotesingle envoyez pas
  de documents confidentiels} sans avoir vérifié les réglages réseau.
\end{itemize}

\begin{center}\rule{0.5\linewidth}{0.5pt}\end{center}

\subsection{3. Recette : transformer un PDF en présentation
PowerPoint}\label{3-recette--transformer-un-pdf-en-pruxe9sentation-powerpoint}

\textbf{Cas d\textquotesingle usage :} vous avez un rapport ou un
mémoire en PDF, et vous voulez un diaporama de présentation. Ici, le
mémoire de CREDIX devient le support de soutenance.

\subsubsection{3.1 Étape par étape}\label{31-uxe9tape-par-uxe9tape}

\begin{enumerate}
\def\labelenumi{\arabic{enumi}.}
\tightlist
\item
  \textbf{Activez} la création de fichiers (section 2.2).
\item
  \textbf{Ouvrez une nouvelle conversation}, de préférence \textbf{dans
  un Projet} (guide 1, section 6), où vous aurez déposé vos consignes.
\item
  \textbf{Joignez le PDF} avec le bouton de pièce jointe. Vérifiez sa
  taille (moins de 30 Mo) et son nombre de pages (section 2.3).
\item
  \textbf{Collez le prompt} de la section 3.2 et complétez les crochets.
\item
  \textbf{Attendez le fichier} \texttt{.pptx} : Claude
  l\textquotesingle affiche dans la conversation, avec un bouton pour le
  télécharger.
\item
  \textbf{Téléchargez-le}, ouvrez-le (PowerPoint ou LibreOffice Impress)
  et \textbf{relisez tout} (section 3.3).
\item
  \textbf{Demandez les corrections} dans la même conversation (section
  3.4).
\end{enumerate}

\CredixCapture{Un PDF joint à une conversation}{../images/guides/04-02_pdf_joint.png}

\emph{Figure 2. Le PDF joint à la conversation.}

\subsubsection{3.2 Le prompt à copier}\label{32-le-prompt-uxe0-copier}

\begin{Verbatim}[breaklines=true,breakanywhere=true,fontsize=\footnotesize,frame=leftline,framerule=1.6pt,rulecolor=\color{credixblue},framesep=3mm,xleftmargin=2mm,xrightmargin=1mm]
Voici mon mémoire en PDF. Crée une présentation PowerPoint (.pptx) de
soutenance à partir de ce document.

Contexte :
- Public : [un jury de trois enseignants, non spécialistes du sujet]
- Durée : [20] minutes, soit environ [15] diapositives
- Langue : français

Structure demandée :
1. Titre et présentation
2. Contexte et problème posé
3. Objectifs
4. Méthode
5. Résultats principaux
6. Limites
7. Conclusion et perspectives
8. Questions

Règles à respecter :
- Une idée par diapositive, un titre clair, cinq lignes au plus.
- Aucun chiffre inventé : tout chiffre doit venir du PDF. Pour chacun,
  indique la page du PDF dans les notes de la diapositive.
- Ajoute des notes de présentation (ce que je dirai à l'oral).
- Style sobre : fond clair, une couleur principale, texte lisible de loin.
- À la fin, donne-moi la liste des diapositives avec leur titre, et
  signale tout passage que tu n'as pas pu résumer fidèlement.
\end{Verbatim}

\subsubsection{3.3 Ce qu\textquotesingle il faut
vérifier}\label{33-ce-quil-faut-vuxe9rifier}

Claude peut se tromper, y compris avec assurance. Avant
d\textquotesingle utiliser le diaporama :

\begin{longtable}[]{@{}
  >{\raggedright\arraybackslash}p{(\linewidth - 2\tabcolsep) * \real{0.2627}}
  >{\raggedright\arraybackslash}p{(\linewidth - 2\tabcolsep) * \real{0.7073}}@{}}
\toprule\noalign{}
\begin{minipage}[b]{\linewidth}\raggedright
Vérification
\end{minipage} & \begin{minipage}[b]{\linewidth}\raggedright
Comment
\end{minipage} \\
\midrule\noalign{}
\endhead
\bottomrule\noalign{}
\endlastfoot
\textbf{Les chiffres} & Comparez chaque chiffre de chaque diapositive
avec le PDF (aide : les numéros de page dans les notes) \\
\textbf{Le nombre de diapositives} & Correspond-il à la durée ? Comptez
environ une diapositive par minute \\
\textbf{Les graphiques et schémas} & Sont-ils repris, et lisibles ?
Sinon, insérez vos images vous-même \\
\textbf{Les mots} & Les noms de méthodes, de variables,
d\textquotesingle auteurs sont-ils correctement écrits ? \\
\textbf{La cohérence} & Le diaporama dit-il la même chose que le mémoire
? Rien d\textquotesingle ajouté, rien de déformé ? \\
\end{longtable}

\CredixCapture{Le diaporama créé, ouvert dans PowerPoint ou LibreOffice}{../images/guides/04-03_resultat_pptx.png}

\emph{Figure 3. Le fichier \texttt{.pptx} produit, ouvert pour
relecture.}

\subsubsection{3.4 Demander des
corrections}\label{34-demander-des-corrections}

Restez dans la \textbf{même conversation} : Claude garde le fichier en
mémoire et le modifie.

\begin{Verbatim}[breaklines=true,breakanywhere=true,fontsize=\footnotesize,frame=leftline,framerule=1.6pt,rulecolor=\color{credixblue},framesep=3mm,xleftmargin=2mm,xrightmargin=1mm]
Réduis la présentation à 12 diapositives : fusionne le contexte et
les objectifs, et supprime la diapositive sur les limites de la méthode.
\end{Verbatim}

\begin{Verbatim}[breaklines=true,breakanywhere=true,fontsize=\footnotesize,frame=leftline,framerule=1.6pt,rulecolor=\color{credixblue},framesep=3mm,xleftmargin=2mm,xrightmargin=1mm]
La diapositive 6 est trop chargée : sépare-la en deux, avec un
graphique par diapositive.
\end{Verbatim}

\begin{Verbatim}[breaklines=true,breakanywhere=true,fontsize=\footnotesize,frame=leftline,framerule=1.6pt,rulecolor=\color{credixblue},framesep=3mm,xleftmargin=2mm,xrightmargin=1mm]
Les notes de la diapositive 4 sont trop longues : garde trois phrases.
\end{Verbatim}

\subsubsection{3.5 Une variante sans intelligence
artificielle}\label{35-une-variante-sans-intelligence-artificielle}

Si vous avez une version \textbf{Markdown} de votre document (comme les
guides de ce dossier), la commande \texttt{pandoc} fabrique un
\texttt{.pptx} \textbf{sans rien changer au texte} (section 8.4). Le
résultat est plus austère, mais \textbf{fidèle au mot près}.

\begin{center}\rule{0.5\linewidth}{0.5pt}\end{center}

\subsection{4. D\textquotesingle autres recettes du même
genre}\label{4-dautres-recettes-du-muxeame-genre}

\begin{longtable}[]{@{}
  >{\raggedright\arraybackslash}p{(\linewidth - 6\tabcolsep) * \real{0.2000}}
  >{\raggedright\arraybackslash}p{(\linewidth - 6\tabcolsep) * \real{0.1995}}
  >{\raggedright\arraybackslash}p{(\linewidth - 6\tabcolsep) * \real{0.3580}}
  >{\raggedright\arraybackslash}p{(\linewidth - 6\tabcolsep) * \real{0.2125}}@{}}
\toprule\noalign{}
\begin{minipage}[b]{\linewidth}\raggedright
Besoin
\end{minipage} & \begin{minipage}[b]{\linewidth}\raggedright
Fichier joint
\end{minipage} & \begin{minipage}[b]{\linewidth}\raggedright
Ce que vous demandez
\end{minipage} & \begin{minipage}[b]{\linewidth}\raggedright
Résultat
\end{minipage} \\
\midrule\noalign{}
\endhead
\bottomrule\noalign{}
\endlastfoot
\textbf{Mettre des données en tableau} & Un CSV ou un PDF contenant des
tableaux & « Extrais les tableaux de ce PDF dans un fichier Excel, un
onglet par tableau, et indique les cellules que tu as dû deviner. » &
\texttt{.xlsx} \\
\textbf{Analyser des résultats} & Un CSV de résultats & « Calcule les
moyennes par groupe, trace un histogramme, et donne-moi le graphique en
image. » & \texttt{.png} \\
\textbf{Rédiger un compte rendu} & Vos notes de réunion en texte & «
Rédige un compte rendu Word : décisions, actions, responsables,
échéances. » & \texttt{.docx} \\
\textbf{Résumer un article scientifique} & L\textquotesingle article en
PDF & « Résume en 10 lignes, puis dis en quoi il concerne le choix
d\textquotesingle une méthode de sélection de variables. » & Texte dans
la conversation \\
\textbf{Comparer deux documents} & Deux PDF & « Liste les différences de
contenu entre ces deux versions, section par section. » & Tableau dans
la conversation, ou \texttt{.xlsx} \\
\textbf{Corriger un document Word} & Un \texttt{.docx} & « Corrige
l\textquotesingle orthographe et les tournures, en gardant ma mise en
forme. Liste tes modifications. » & \texttt{.docx} corrigé \\
\end{longtable}

\textbf{Règle commune :} demandez toujours à Claude de \textbf{signaler
ce qu\textquotesingle il a deviné ou n\textquotesingle a pas pu lire}.
C\textquotesingle est là que se cachent les erreurs.

\begin{center}\rule{0.5\linewidth}{0.5pt}\end{center}

\subsection{5. Les Skills : apprendre à Claude votre façon de
travailler}\label{5-les-skills--apprendre-uxe0-claude-votre-fauxe7on-de-travailler}

\subsubsection{5.1 Ce que c\textquotesingle est}\label{51-ce-que-cest}

Un \textbf{Skill} (« compétence ») est un \textbf{dossier
d\textquotesingle instructions}, parfois accompagné de scripts, que
Claude charge \textbf{uniquement quand la tâche le demande}. Il enseigne
à Claude \emph{comment} faire une tâche précise, de façon répétable.

\begin{longtable}[]{@{}
  >{\raggedright\arraybackslash}p{(\linewidth - 4\tabcolsep) * \real{0.1552}}
  >{\raggedright\arraybackslash}p{(\linewidth - 4\tabcolsep) * \real{0.3259}}
  >{\raggedright\arraybackslash}p{(\linewidth - 4\tabcolsep) * \real{0.4889}}@{}}
\toprule\noalign{}
\begin{minipage}[b]{\linewidth}\raggedright
\end{minipage} & \begin{minipage}[b]{\linewidth}\raggedright
Projet
\end{minipage} & \begin{minipage}[b]{\linewidth}\raggedright
Skill
\end{minipage} \\
\midrule\noalign{}
\endhead
\bottomrule\noalign{}
\endlastfoot
\textbf{Nature} & Documents et consignes générales & Une
\textbf{procédure} pour une tâche \\
\textbf{Quand c\textquotesingle est utilisé} & Toujours, dans les
conversations du projet & \textbf{Seulement} si la demande correspond \\
\textbf{Exemple} & Le cahier des charges de CREDIX & « Comment rédiger
un diagramme UML conforme à nos règles » \\
\end{longtable}

Il en existe de plusieurs sortes :

\begin{itemize}
\tightlist
\item
  les Skills \textbf{d\textquotesingle Anthropic} (création de fichiers
  Excel, Word, PowerPoint, PDF), que Claude utilise
  \textbf{automatiquement} ;
\item
  les Skills \textbf{personnels}, que vous écrivez (en Markdown : aucune
  programmation nécessaire pour un Skill simple) ;
\item
  les Skills \textbf{d\textquotesingle organisation}, distribués par les
  administrateurs d\textquotesingle une équipe ;
\item
  des Skills de \textbf{partenaires} (Notion, Figma, Atlassian, etc.).
\end{itemize}

Les Skills sont disponibles sur tous les abonnements, et demandent que
l\textquotesingle option de création de fichiers soit activée (section
2.2).

\subsubsection{5.2 Où les trouver dans
claude.ai}\label{52-ouxf9-les-trouver-dans-claudeai}

\begin{enumerate}
\def\labelenumi{\arabic{enumi}.}
\tightlist
\item
  Cliquez sur \textbf{Personnaliser} (\emph{Customize}) dans votre
  compte.
\item
  Ouvrez \textbf{Skills}.
\item
  Cliquez sur \textbf{+}, puis \textbf{Parcourir les skills}
  (\emph{Browse skills}) pour ouvrir l\textquotesingle annuaire.
\end{enumerate}

\CredixCapture{Menu Personnaliser, liste des Skills}{../images/guides/04-04_skills_annuaire.png}

\emph{Figure 4. L\textquotesingle annuaire des Skills.}

\subsubsection{5.3 Dans Claude Code : écrire son propre
Skill}\label{53-dans-claude-code--uxe9crire-son-propre-skill}

Un Skill est un dossier contenant un fichier \textbf{\texttt{SKILL.md}},
qui commence par un en-tête (nom et description) puis contient les
instructions.

\begin{longtable}[]{@{}
  >{\raggedright\arraybackslash}p{(\linewidth - 2\tabcolsep) * \real{0.4157}}
  >{\raggedright\arraybackslash}p{(\linewidth - 2\tabcolsep) * \real{0.5543}}@{}}
\toprule\noalign{}
\begin{minipage}[b]{\linewidth}\raggedright
Portée
\end{minipage} & \begin{minipage}[b]{\linewidth}\raggedright
Emplacement
\end{minipage} \\
\midrule\noalign{}
\endhead
\bottomrule\noalign{}
\endlastfoot
\textbf{Personnel} (tous vos projets) &
\texttt{\textasciitilde{}/.claude/skills/\textless{}nom-du-skill\textgreater{}/SKILL.md} \\
\textbf{Projet} (ce dépôt uniquement) &
\texttt{.claude/skills/\textless{}nom-du-skill\textgreater{}/SKILL.md} \\
\end{longtable}

Exemple minimal, à placer dans
\texttt{.\allowbreak{}claude/\allowbreak{}skills/\allowbreak{}relecture-\allowbreak{}conception/\allowbreak{}SKILL.\allowbreak{}md}
:

\begin{Verbatim}[breaklines=true,breakanywhere=true,fontsize=\footnotesize,frame=leftline,framerule=1.6pt,rulecolor=\color{credixblue},framesep=3mm,xleftmargin=2mm,xrightmargin=1mm]
---
name: relecture-conception
description: Relit un document de conception et signale les exigences sans critère de test, les ambiguïtés et les incohérences de vocabulaire.
---

Quand on te demande de relire un document de conception :

1. Liste chaque exigence et indique si elle a un critère de test mesurable.
2. Signale les termes utilisés avec deux sens différents.
3. Ne réécris rien : propose des corrections sous forme de liste.
4. Termine par les cinq problèmes les plus importants.
\end{Verbatim}

Le Skill se déclenche \textbf{de deux façons} :

\begin{itemize}
\tightlist
\item
  \textbf{automatiquement}, quand votre demande correspond à sa
  description (d\textquotesingle où l\textquotesingle importance
  d\textquotesingle une description précise) ;
\item
  \textbf{manuellement}, en tapant
  \texttt{/\allowbreak{}relecture-\allowbreak{}conception} dans Claude
  Code.
\end{itemize}

\begin{quote}
\textbf{Un exemple concret dans ce projet.} Un Skill « génie logiciel »
regroupe les règles de rédaction des diagrammes UML (une association
plutôt qu\textquotesingle un attribut, multiplicités aux deux
extrémités, etc.). Il a servi à écrire les diagrammes du guide 3, sans
qu\textquotesingle on ait à répéter ces règles à chaque conversation.
\end{quote}

\begin{center}\rule{0.5\linewidth}{0.5pt}\end{center}

\subsection{6. Les artefacts et les
connecteurs}\label{6-les-artefacts-et-les-connecteurs}

\subsubsection{6.1 Les artefacts}\label{61-les-artefacts}

Un \textbf{artefact} est un contenu que Claude produit et qui
\textbf{s\textquotesingle ouvre à côté de la conversation} : vous pouvez
le modifier, y revenir et le partager par un lien.

Il peut s\textquotesingle agir de : un \textbf{document} (Markdown ou
texte), un \textbf{extrait de code}, une \textbf{page web}
d\textquotesingle une seule page (HTML), une image \textbf{SVG}, un
\textbf{schéma} ou un organigramme, un \textbf{composant interactif}.

\begin{itemize}
\tightlist
\item
  Pour modifier un artefact, demandez à Claude, ou, pour un document
  Markdown, surlignez le passage, cliquez sur \textbf{Edit with Claude}
  et écrivez votre demande.
\item
  En bas à droite de la fenêtre de l\textquotesingle artefact, vous
  pouvez \textbf{voir le code}, \textbf{copier} le contenu ou
  \textbf{télécharger} le fichier.
\item
  La fonction \textbf{Publier et partager} permet de partager un
  artefact par un lien.
\end{itemize}

Les artefacts sont disponibles sur tous les abonnements. Des fonctions
de présentation (Claude Slides) et de documents sont proposées,
\textbf{en version bêta}, sur les abonnements payants.

\CredixCapture{Un artefact ouvert à côté de la conversation}{../images/guides/04-05_artefact.png}

\emph{Figure 5. Un artefact à côté de la conversation.}

\subsubsection{6.2 Les connecteurs}\label{62-les-connecteurs}

Un \textbf{connecteur} relie Claude à un \textbf{service extérieur} pour
qu\textquotesingle il lise ou utilise vos données : Google Workspace
(Gmail, Drive, Agenda), GitHub, Microsoft 365, etc. On peut aussi
ajouter des connecteurs personnalisés.

\textbf{Activation (abonnements Gratuit, Pro et Max) :}

\begin{enumerate}
\def\labelenumi{\arabic{enumi}.}
\tightlist
\item
  Ouvrez \textbf{Personnaliser}, puis \textbf{Connecteurs}
  (\emph{Customize \textgreater{} Connectors}).
\item
  Ajoutez le connecteur voulu et autorisez l\textquotesingle accès quand
  la page vous le demande.
\item
  Dans une conversation, cliquez sur le bouton \textbf{+}, puis
  \textbf{Connecteurs}, pour l\textquotesingle activer pour cette
  conversation.
\end{enumerate}

Sur les abonnements Team et Enterprise, un propriétaire ajoute
d\textquotesingle abord les connecteurs dans \textbf{Paramètres de
l\textquotesingle organisation}.

\CredixCapture{Menu Connecteurs}{../images/guides/04-06_connecteurs.png}

\emph{Figure 6. Le menu des connecteurs.}

\begin{quote}
\textbf{Prudence.} Ne connectez Claude qu\textquotesingle à des services
\textbf{de confiance}. Un connecteur peut \textbf{lire, créer, modifier
ou supprimer} des données. Lisez attentivement les autorisations
demandées, et n\textquotesingle acceptez que celles dont vous avez
besoin.
\end{quote}

\begin{center}\rule{0.5\linewidth}{0.5pt}\end{center}

\subsection{7. Claude Code : au-delà de l\textquotesingle écriture de
code}\label{7-claude-code--au-deluxe0-de-luxe9criture-de-code}

Claude Code (guide 1, sections 7 et 8) sait faire plus que programmer.
D\textquotesingle après la documentation officielle :

\begin{longtable}[]{@{}
  >{\raggedright\arraybackslash}p{(\linewidth - 2\tabcolsep) * \real{0.2071}}
  >{\raggedright\arraybackslash}p{(\linewidth - 2\tabcolsep) * \real{0.7629}}@{}}
\toprule\noalign{}
\begin{minipage}[b]{\linewidth}\raggedright
Fonction
\end{minipage} & \begin{minipage}[b]{\linewidth}\raggedright
À quoi elle sert
\end{minipage} \\
\midrule\noalign{}
\endhead
\bottomrule\noalign{}
\endlastfoot
\textbf{\texttt{CLAUDE.md}} & Fichier lu au début de chaque session :
règles du projet, commandes, conventions \\
\textbf{Skills} & Procédures réutilisables (\texttt{/nom-du-skill}),
partageables avec l\textquotesingle équipe (section 5.3) \\
\textbf{Hooks} & Commandes lancées automatiquement avant ou après une
action de Claude (par exemple : reformater un fichier après chaque
modification) \\
\textbf{Sous-agents} & Plusieurs agents Claude qui travaillent en
parallèle sur des parties différentes d\textquotesingle une tâche \\
\textbf{MCP} & Connexion à des outils extérieurs (documents Google
Drive, tickets Jira, Slack\ldots) \\
\textbf{Git} & Enregistrer des modifications, créer des branches,
rédiger les messages, ouvrir des demandes de fusion \\
\textbf{Mode non interactif} & Lancer une seule demande depuis le
terminal, par exemple \texttt{claude\ -p\ "ta\ demande"}, pour
l\textquotesingle enchaîner avec d\textquotesingle autres commandes \\
\textbf{Tâches planifiées} & Répéter une demande à intervalles
réguliers \\
\textbf{Autres surfaces} & Terminal, VS Code, JetBrains, application de
bureau, navigateur (claude.ai/code) \\
\end{longtable}

\textbf{Une idée utile pour ce projet :} vous pouvez demander à Claude
Code de \textbf{transformer un document} (par exemple, écrire un fichier
\texttt{.pptx} ou \texttt{.docx} à partir d\textquotesingle un
\texttt{.md}), puis de \textbf{relire ce qu\textquotesingle il a
produit}. Comme pour tout ce qui précède, \textbf{relisez le résultat}.

\begin{center}\rule{0.5\linewidth}{0.5pt}\end{center}

\subsection{8. Fabriquer soi-même des PDF, Word et PowerPoint à partir
de
Markdown}\label{8-fabriquer-soi-muxeame-des-pdf-word-et-powerpoint-uxe0-partir-de-markdown}

Tous les PDF de ce dossier sont fabriqués \textbf{sans intelligence
artificielle}, à partir des fichiers Markdown, avec deux outils libres :
\textbf{pandoc} (convertisseur de documents) et \textbf{LaTeX} (mise en
page). Avantage : le texte est identique au mot près, et vous refaites
le PDF en une commande dès que le \texttt{.md} change.

\subsubsection{8.1 Installer les outils
(Linux)}\label{81-installer-les-outils-linux}

\begin{Verbatim}[breaklines=true,breakanywhere=true,fontsize=\footnotesize,frame=leftline,framerule=1.6pt,rulecolor=\color{credixblue},framesep=3mm,xleftmargin=2mm,xrightmargin=1mm]
sudo apt update
sudo apt install pandoc texlive-latex-recommended texlive-latex-extra \
  texlive-fonts-recommended texlive-lang-french poppler-utils
\end{Verbatim}

Vérifiez :

\begin{Verbatim}[breaklines=true,breakanywhere=true,fontsize=\footnotesize,frame=leftline,framerule=1.6pt,rulecolor=\color{credixblue},framesep=3mm,xleftmargin=2mm,xrightmargin=1mm]
pandoc --version | head -1
pdflatex --version | head -1
\end{Verbatim}

Chacune doit afficher un numéro de version. Le paquet
\texttt{poppler-utils} fournit \texttt{pdfunite} (assembler des PDF) et
\texttt{pdfinfo} (compter les pages).

\subsubsection{8.2 Fabriquer un PDF et son fichier
LaTeX}\label{82-fabriquer-un-pdf-et-son-fichier-latex}

Depuis la racine du dépôt :

\begin{Verbatim}[breaklines=true,breakanywhere=true,fontsize=\footnotesize,frame=leftline,framerule=1.6pt,rulecolor=\color{credixblue},framesep=3mm,xleftmargin=2mm,xrightmargin=1mm]
docs/build/build_pdf.sh docs/guides/04_autres_fonctions_claude.md \
  /tmp/exemple.pdf "Guide 4" "Autres fonctions de Claude"
\end{Verbatim}

Les quatre paramètres sont : le fichier Markdown, le PDF à produire, le
titre de la page de garde, son sous-titre. Le script fait deux choses :

\begin{enumerate}
\def\labelenumi{\arabic{enumi}.}
\tightlist
\item
  il convertit le Markdown en un fichier \textbf{LaTeX autonome}
  (\texttt{/\allowbreak{}tmp/\allowbreak{}exemple.\allowbreak{}tex}),
  avec la page de garde, la table des matières et les en-têtes ;
\item
  il \textbf{compile ce fichier} en PDF
  (\texttt{/\allowbreak{}tmp/\allowbreak{}exemple.\allowbreak{}pdf}).
\end{enumerate}

Pour tout le dossier de remise, une seule commande refait les six PDF,
leurs fichiers \texttt{.tex} (dans \texttt{docs/tex/}) et le PDF
regroupé :

\begin{Verbatim}[breaklines=true,breakanywhere=true,fontsize=\footnotesize,frame=leftline,framerule=1.6pt,rulecolor=\color{credixblue},framesep=3mm,xleftmargin=2mm,xrightmargin=1mm]
docs/build/build_all.sh
\end{Verbatim}

Pour compter les pages d\textquotesingle un PDF :

\begin{Verbatim}[breaklines=true,breakanywhere=true,fontsize=\footnotesize,frame=leftline,framerule=1.6pt,rulecolor=\color{credixblue},framesep=3mm,xleftmargin=2mm,xrightmargin=1mm]
pdfinfo /tmp/exemple.pdf | grep Pages
\end{Verbatim}

\subsubsection{\texorpdfstring{8.3 Insérer des images dans le fichier
\texttt{.tex}}{8.3 Insérer des images dans le fichier .tex}}\label{83-insuxe9rer-des-images-dans-le-fichier-tex}

Les documents contiennent des cadres \textbf{« CAPTURE À INSÉRER »} aux
endroits où une capture d\textquotesingle écran est attendue. Deux
façons de les remplacer :

\textbf{Façon 1, sans rien modifier (la plus simple).} Le cadre indique
un nom de fichier, par exemple
\texttt{.\allowbreak{}.\allowbreak{}/\allowbreak{}images/\allowbreak{}guides/\allowbreak{}01-\allowbreak{}01\_\allowbreak{}claude\_\allowbreak{}ai\_\allowbreak{}accueil.\allowbreak{}png}.
Déposez votre capture \textbf{sous ce nom exact} dans
\texttt{docs/\allowbreak{}images/\allowbreak{}guides/\allowbreak{}},
puis recompilez, deux fois, depuis \texttt{docs/tex/} :

\begin{Verbatim}[breaklines=true,breakanywhere=true,fontsize=\footnotesize,frame=leftline,framerule=1.6pt,rulecolor=\color{credixblue},framesep=3mm,xleftmargin=2mm,xrightmargin=1mm]
cd docs/tex
pdflatex 02_Guide1_Outils_Claude.tex
pdflatex 02_Guide1_Outils_Claude.tex
\end{Verbatim}

La capture apparaît à la place du cadre. La deuxième compilation met à
jour la table des matières. Le PDF est créé à côté du \texttt{.tex}.

\textbf{Façon 2, en modifiant le \texttt{.tex}.} Cherchez la ligne du
cadre :

\begin{Verbatim}[breaklines=true,breakanywhere=true,fontsize=\footnotesize,frame=leftline,framerule=1.6pt,rulecolor=\color{credixblue},framesep=3mm,xleftmargin=2mm,xrightmargin=1mm]
\CredixCapture{Page d'accueil de claude.ai}{../images/guides/01-01_claude_ai_accueil.png}
\end{Verbatim}

et remplacez-la par votre propre commande :

\begin{Verbatim}[breaklines=true,breakanywhere=true,fontsize=\footnotesize,frame=leftline,framerule=1.6pt,rulecolor=\color{credixblue},framesep=3mm,xleftmargin=2mm,xrightmargin=1mm]
\begin{center}
  \includegraphics[width=\linewidth]{../images/guides/mon_image.png}
\end{center}
\end{Verbatim}

\begin{quote}
\textbf{Attention.} Le \texttt{.tex} est \textbf{fabriqué} à partir du
Markdown : quand on relance
\texttt{docs/\allowbreak{}build/\allowbreak{}build\_\allowbreak{}all.\allowbreak{}sh},
il est \textbf{réécrit}, et les modifications faites directement dedans
sont perdues. Choisissez donc une source : le Markdown (recommandé) ou
le \texttt{.tex}. Le fichier
\texttt{docs/\allowbreak{}tex/\allowbreak{}LISEZMOI.\allowbreak{}md}
résume la marche à suivre.
\end{quote}

\subsubsection{8.4 Fabriquer un fichier Word ou
PowerPoint}\label{84-fabriquer-un-fichier-word-ou-powerpoint}

Depuis un fichier Markdown, \texttt{pandoc} produit directement :

\begin{Verbatim}[breaklines=true,breakanywhere=true,fontsize=\footnotesize,frame=leftline,framerule=1.6pt,rulecolor=\color{credixblue},framesep=3mm,xleftmargin=2mm,xrightmargin=1mm]
pandoc -f gfm mon_document.md -o mon_document.docx
pandoc -f gfm mon_document.md -o ma_presentation.pptx
pandoc -f gfm mon_document.md -s -o ma_page.html
\end{Verbatim}

Vérifié sur un petit document d\textquotesingle essai : les trois
commandes réussissent. Pour la présentation, \textbf{un titre de niveau
1 (\texttt{\#}) donne la diapositive de titre, et chaque titre de niveau
2 (\texttt{\#\#}) donne une diapositive}, dont le contenu est le texte
placé dessous. Structurez donc votre Markdown en conséquence : peu de
texte sous chaque titre \texttt{\#\#}.

\subsubsection{8.5 Quand utiliser quoi ?}\label{85-quand-utiliser-quoi-}

\begin{longtable}[]{@{}
  >{\raggedright\arraybackslash}p{(\linewidth - 2\tabcolsep) * \real{0.6406}}
  >{\raggedright\arraybackslash}p{(\linewidth - 2\tabcolsep) * \real{0.3294}}@{}}
\toprule\noalign{}
\begin{minipage}[b]{\linewidth}\raggedright
Besoin
\end{minipage} & \begin{minipage}[b]{\linewidth}\raggedright
Outil
\end{minipage} \\
\midrule\noalign{}
\endhead
\bottomrule\noalign{}
\endlastfoot
Un diaporama \textbf{rédigé et mis en forme} à partir
d\textquotesingle un PDF, avec résumé et notes & \textbf{Claude}
(section 3) \\
Un diaporama, un Word ou un PDF \textbf{exactement fidèle} au texte
d\textquotesingle un \texttt{.md} & \textbf{pandoc} (section 8) \\
Un PDF soigné (page de garde, table des matières), et le fichier LaTeX
pour y ajouter des images & \textbf{pandoc + LaTeX} (sections 8.2 et
8.3) \\
\end{longtable}

\begin{center}\rule{0.5\linewidth}{0.5pt}\end{center}

\subsection{9. Quel outil pour quel besoin
?}\label{9-quel-outil-pour-quel-besoin-}

\begin{longtable}[]{@{}
  >{\raggedright\arraybackslash}p{(\linewidth - 4\tabcolsep) * \real{0.4128}}
  >{\raggedright\arraybackslash}p{(\linewidth - 4\tabcolsep) * \real{0.3646}}
  >{\raggedright\arraybackslash}p{(\linewidth - 4\tabcolsep) * \real{0.1926}}@{}}
\toprule\noalign{}
\begin{minipage}[b]{\linewidth}\raggedright
Je veux\ldots{}
\end{minipage} & \begin{minipage}[b]{\linewidth}\raggedright
J\textquotesingle utilise
\end{minipage} & \begin{minipage}[b]{\linewidth}\raggedright
Où lire
\end{minipage} \\
\midrule\noalign{}
\endhead
\bottomrule\noalign{}
\endlastfoot
Réfléchir, cadrer un besoin, chercher dans la littérature & claude.ai,
en mode \textbf{Projet} & Guide 1, sections 5 et 6 \\
Concevoir une application, de A à Z & La \textbf{méthode} (Projet,
sessions, documents, décisions) & Guide 2 \\
Écrire, tester et corriger le code & \textbf{Claude Code}, en mode
\textbf{Plan} & Guide 1, section 8 ; guide 3 \\
Créer un Excel, un Word ou un PowerPoint & claude.ai avec
\textbf{création de fichiers} & Sections 2 à 4 \\
Que Claude applique toujours mes règles pour une tâche & Un
\textbf{Skill} & Section 5 \\
Que Claude lise mes documents Drive ou mon dépôt GitHub & Un
\textbf{connecteur} & Section 6.2 \\
Un PDF ou un document fidèle au mot près & \textbf{pandoc} & Section
8 \\
\end{longtable}

\begin{center}\rule{0.5\linewidth}{0.5pt}\end{center}

\subsection{10. Précautions}\label{10-pruxe9cautions}

\begin{itemize}
\tightlist
\item
  \textbf{Relisez tout.} Un fichier bien présenté n\textquotesingle est
  pas forcément exact. Chiffres, noms et références se vérifient à la
  main.
\item
  \textbf{Confidentialité.} Ne donnez à Claude ni mot de passe, ni clé
  d\textquotesingle API, ni donnée personnelle de clients réels.
  Vérifiez les réglages réseau si des fichiers sont créés (section 2.4).
\item
  \textbf{Connecteurs et Skills de tiers.} N\textquotesingle installez
  que ceux dont la source est fiable : un Skill peut contenir des
  scripts, un connecteur peut modifier vos données.
\item
  \textbf{Gardez la source.} Conservez toujours
  l\textquotesingle original (le PDF, le Markdown) : le fichier produit
  par Claude en est une \textbf{dérivation}, pas le remplacement.
\item
  \textbf{Fonctions en évolution.} Les menus, les limites et les
  fonctions en version bêta changent : en cas de doute, la page
  officielle fait foi.
\end{itemize}

\begin{center}\rule{0.5\linewidth}{0.5pt}\end{center}

\subsection{11. Sources officielles}\label{11-sources-officielles}

Pages consultées le 19 septembre 2026 :

\begin{itemize}
\tightlist
\item
  Créer et modifier des fichiers avec Claude :
  \url{https://support.claude.com/en/articles/12111783-create-and-edit-files-with-claude}
\item
  Envoyer des fichiers à Claude (limites, PDF) :
  \url{https://support.claude.com/en/articles/8241126-upload-files-to-claude}
\item
  Qu\textquotesingle est-ce qu\textquotesingle un Skill :
  \url{https://support.claude.com/en/articles/12512176-what-are-skills}
\item
  Skills dans Claude Code : \url{https://code.claude.com/docs/en/skills}
\item
  Artefacts :
  \url{https://support.claude.com/en/articles/9487310-what-are-artifacts-and-how-do-i-use-them}
\item
  Connecteurs personnalisés :
  \url{https://support.claude.com/en/articles/11175166-get-started-with-custom-connectors-using-remote-mcp}
\item
  Claude Code, vue d\textquotesingle ensemble :
  \url{https://code.claude.com/docs/en/overview}
\end{itemize}

\end{document}
